\chapter{Literature Review}
\label{ch:literature}

This chapter reviews four lines of work that meet in this thesis: (1) learned priors for scientific inverse problems,
(2) explicit Gaussian representations, (3) Gaussian and neural-field methods for scientific and medical imaging, and
(4) dynamic representations and sequence models. It closes with the gap this thesis addresses.

\section{Priors for scientific inverse problems}
Classical solvers combine a data-fidelity term with hand-crafted regularisers such as total
variation~\citep{rudin1992tv}. Deep learning added two families of priors. \emph{Untrained} priors use the structure
of a network as the regulariser: Deep Image Prior~\citep{ulyanov2018dip} fits a randomly initialised convolutional
network to a single corrupted image and stops early. \emph{Learned} priors are trained on many examples; diffusion
models are currently the most general, and Diffusion Posterior Sampling~\citep{chung2023dps} guides the reverse
diffusion with the gradient of the data likelihood, so the same prior can serve many forward models.

InverseBench~\citep{zheng2025inversebench} evaluates 14 such plug-and-play diffusion-prior algorithms on five
scientific problems: linear inverse scattering, compressed-sensing MRI on fastMRI~\citep{zbontar2018fastmri}, black
hole imaging, full waveform inversion on OpenFWI~\citep{deng2022openfwi}, and the 2D Navier--Stokes equations. Its
findings that motivate this thesis are that (i) scientific forward models differ strongly from image-restoration
operators (complex-valued, non-linear, expensive, or without gradients), (ii) diffusion priors help most when
measurements are sparse, and (iii) they can fail when outputs violate physical constraints. All its methods represent
the unknown on a pixel grid. Physics-informed neural networks~\citep{raissi2019pinn} and neural operators such as
FNO~\citep{li2021fno} are complementary: they learn solution operators of PDEs rather than priors over unknowns.

\section{Explicit Gaussian representations}
3DGS~\citep{kerbl20233dgs} showed that an explicit, adaptive set of anisotropic Gaussians can match or exceed neural
fields~\citep{mildenhall2020nerf,mueller2022instantngp} in quality while rendering in real time. Its key ingredients
are a differentiable rasteriser, a covariance parametrisation by rotation and scale that stays positive
semi-definite, and adaptive density control. GaussianImage~\citep{zhang2024gaussianimage} transfers the idea to 2D:
an image is represented by 2D Gaussians and rendered by accumulated summation, achieving fast decoding and
compression. Our 2D and 3D representations follow the same principle, but are evaluated on a regular grid and
supervised through a physics operator instead of pixels.

\section{Gaussians and neural fields in scientific imaging}
\paragraph{Tomography.} X-Gaussian~\citep{cai2024xgaussian} adapts Gaussian splatting to X-ray projection imaging
with a radiative Gaussian model, and R$^2$-Gaussian~\citep{zha2024r2gaussian} rectifies the integration bias of
splatting for tomographic volume reconstruction, showing that Gaussians can serve as a voxel-free representation of
attenuation volumes. Both rely on ray-based projection, the natural analogue of camera rendering; scattering
measurements are instead a wave-based integral over the whole object.
\paragraph{Astronomy.} BH-NeRF~\citep{levis2022bhnerf} reconstructs 3D emission around a black hole over time with a
neural field constrained by gravitational lensing and orbital dynamics, an example of physics replacing multi-view
geometry.
\paragraph{Microscopy data.} Simulated cell images with exact ground truth are produced by CytoPacq and related
generators~\citep{wiesner2019cytopacq}; the Broad Bioimage Benchmark Collection~\citep{ljosa2012bbbc} distributes such
synthetic 3D nucleus datasets (BBBC024), which we use as real 3D cell shapes.

\section{Dynamic representations and sequence models}
D-NeRF~\citep{pumarola2021dnerf} and 4D Gaussian Splatting~\citep{wu20244dgs} model motion with a canonical
representation and a deformation field queried independently at each time; HexPlane~\citep{cao2023hexplane} and
K-Planes~\citep{fridovich2023kplanes} factorise space--time features into planes. Dynamic 3D
Gaussians~\citep{luiten2024dynamic} track persistent Gaussians frame by frame, and native 4D
Gaussians~\citep{yang2024realtime4dgs} slice 4D primitives in time. Language-embedded variants
(LangSplat~\citep{qin2024langsplat}, 4D LangSplat~\citep{li2025langsplat4d}) attach semantic features to such
Gaussians, the starting point of this thesis' original direction. None of these treats a Gaussian's trajectory as a
sequence. State-space models~\citep{gu2022s4,gu2023mamba} are efficient sequence models that could do exactly that,
and have been adopted in vision~\citep{zhu2024visionmamba}; their use as a motion model for 4D Gaussians is, to our
knowledge, unexplored.

\section{Synthesis and research gap}
Table~\ref{tab:lit} summarises the landscape. Diffusion-prior methods are general but represent the unknown on a
grid and need large training sets; Gaussian methods are compact and adaptive but have been demonstrated mainly with
camera- or ray-based measurements. The gap addressed here is the use of 3D/4D Gaussian representations as the unknown
in \emph{wave-based and PDE-based scientific inverse problems}, the analysis of their failure modes under severe
ill-posedness, and sequence-model motion priors for their dynamics.

\begin{table}[t]
  \centering
  \small
  \caption{Representative related work. ``Meas.'' is the kind of measurement used for supervision.}
  \label{tab:lit}
  \begin{tabular}{p{3.6cm}p{2.6cm}p{2.7cm}p{4.6cm}}
    \toprule
    Work & Representation & Meas. & Relation to this thesis \\
    \midrule
    NeRF~\citep{mildenhall2020nerf} & MLP field & RGB images & implicit field; rendering by rays \\
    3DGS~\citep{kerbl20233dgs} & 3D Gaussians & RGB images & our representation, without cameras \\
    4DGS~\citep{wu20244dgs} & Gaussians + deformation & video & our 4D design (MLP deformation) \\
    GaussianImage~\citep{zhang2024gaussianimage} & 2D Gaussians & pixels & 2D Gaussian rendering on a grid \\
    X-Gaussian~\citep{cai2024xgaussian}, R$^2$-Gaussian~\citep{zha2024r2gaussian} & 3D Gaussians & X-ray projections & Gaussians for physics, ray-based \\
    BH-NeRF~\citep{levis2022bhnerf} & neural field & radio interferometry & physics replaces multi-view geometry \\
    DPS~\citep{chung2023dps}, InverseBench~\citep{zheng2025inversebench} & pixel grid + diffusion prior & physics operators & same problems, learned prior \\
    Mamba~\citep{gu2023mamba} & selective SSM & sequences & our motion model over time \\
    \textbf{This thesis} & 2D/3D/4D Gaussians & scattering; sparse flow samples & Gaussians + physics; Mamba motion \\
    \bottomrule
  \end{tabular}
\end{table}
